\chapter{Literature Review}
\label{ch:litreview}

This thesis is situated within the intersection of three distinct research fields: computational bus service design and optimization, transportation equity measurement and integration, and applied category theory for complex systems engineering. I propose a general-purpose framework for bus network redesigns that natively incorporates equity and accessibility alongside service cost and quality, using the monotone theory of co-design within applied category theory. Accordingly, this review synthesizes existing scholarship on public transit network design and routing optimization, the evolving methods for defining and measuring equity in transportation, the emerging application of categorical and compositional mathematical frameworks to mobility systems, and the data infrastructure (including \gls{gtfs}, \gls{osm}, and Census/\gls{acs} data) that enables data-integrated transit planning. The review also examines school bus transit optimization as the case study domain, drawing on the extensive literature on the \gls{sbrp}.


\section{Bus Network Design and Optimization}

\subsection{Foundational Approaches}

The problem of designing efficient bus networks has been a cornerstone of transportation operations research for decades. \textcite{cederBusNetworkDesign1986} proposed one of the earliest systematic approaches, decomposing bus network design into three sequential stages: trip distribution analysis, route design, and frequency setting. Their work recognized that route planning decisions had received far less attention than lower-level problems (such as vehicle and driver scheduling), and argued that periodic redesign of bus networks could significantly improve system performance. This hierarchical decomposition has since become a standard paradigm in the field.

\textcite{furthSettingFrequenciesBus1981} built on this foundation by developing a model for setting frequencies on bus routes that maximizes net social benefit subject to constraints on total subsidy, fleet size, and maximum headways. This formulation explicitly modeled two objectives, including consumer surplus from reduced waiting time and social ridership benefits from externalities, and demonstrated that optimal allocation could shift significant resources from peak to midday periods through a case study of the \gls{mbta} system. The model's finding that the relative weights given to wait-time savings versus ridership benefits had little impact on optimal allocation was a notable result, suggesting robustness in the underlying economic structure of transit frequency-setting.

\subsection{Integrated Optimization in Public Transport Planning}

A central challenge in public transport planning is the integration of multiple planning stages, including line planning, timetabling, and vehicle scheduling. These problems have traditionally been solved sequentially but are inherently interdependent. Thus, \textcite{schieweIntegratedOptimizationPublic2020} provided a comprehensive treatment of integrated optimization in public transport planning, developing both exact and heuristic approaches for combining these stages. This iterative re-optimization approach allows modifications to an existing public transport plan by sequentially applying algorithms for line planning, timetabling, and vehicle scheduling, with theoretical convergence guarantees under certain cost assumptions. The work demonstrates that iterating between these stages monotonically decreases travel time when re-timetabling and re-vehicle-scheduling are applied, establishing formal convergence properties via the monotone convergence theorem.

Going further, \textcite{patzoldLookAheadApproachesIntegrated2017} extended the integration paradigm by developing look-ahead approaches that incorporate information from downstream planning stages into upstream decisions. Their work addresses the fundamental tension in sequential planning: decisions made early in the process may end up being detrimental when evaluated against the full set of objectives. By enabling line planning to ``look ahead'' to timetabling and vehicle scheduling implications, their framework produces solutions that better balance passenger- and cost-oriented objectives.

\section{School Bus Routing and Scheduling}
\subsection{The School Bus Routing Problem}

The \gls{sbrp}, a specialized variant of the \gls{vrp}, has attracted substantial and accelerating research interest since the 1970s. \textcite{wang2025School} provide a comprehensive review of 127 peer-reviewed studies published between 1970 and 2024, proposing a classification framework based on problem characteristics such as the number of schools served, service environment, fleet composition, and operational constraints. Their analysis reveals an evolution from early cost-minimization models toward more sophisticated multi-objective frameworks that incorporate heterogeneous fleets, service equity, and dynamic constraints.

In many of these studies, the \gls{sbrp} is typically decomposed into several connected sub-problems: data preparation, stop selection, route generation, route scheduling, and school timetabling. Route generation has consistently attracted the largest share of scholarly attention, while school timetabling remains comparatively underexplored. Conventional \gls{sbrp} models are typically formulated as \glspl{milp}, with objectives including minimizing the number of buses, total travel distance, and total student riding time, subject to constraints on vehicle capacity, maximum riding time, time windows, and walking distances. A notable recent trend is the inclusion of carbon emission constraints, signaling a growing emphasis on environmental sustainability.

In terms of solution approaches, metaheuristic algorithms such as genetic algorithms, simulated annealing, and adaptive large neighborhood search have become pervasive for addressing the $\mathcal{NP}$-hard complexity of large-scale \gls{sbrp} variants. The field is also seeing emerging integration of artificial intelligence techniques, including reinforcement learning for dynamic dispatching and deep reinforcement learning for end-to-end routing policies.

\subsection{Optimization-Driven Policy: The Boston Case}

The work conducted by \textcite{bertsimasOptimizingSchoolsStart2019} represents a landmark application of operations research to school transportation policy. They developed the \gls{bird} algorithm, which bridges the gap between single-school and multi-school routing subproblems by generating multiple Pareto-optimal routing scenarios for each school, then jointly selecting the best combination across the district. \Gls{bird} significantly outperformed state-of-the-art routing methods, achieving improvements of 4 to 12\% on benchmark datasets.

As noted in \textcite{bertsimasBusRoutingOptimization2020}, \gls{bird}'s implementation in \gls{bps} led to a reduction of 50 buses from the school district's fleet, generating approximately \$5 million in annual savings while maintaining service quality for students. Critically, the work extended beyond routing to address the \gls{stsp}, formulated as a multiobjective \gls{gqap} that captures pairwise routing compatibility costs between schools. The \gls{stsp} formulation enabled \gls{bps} to explore tradeoffs between competing priorities (maximizing the number of high school students starting after 8:00 AM, minimizing elementary students ending after 4:00 PM, reinvesting transportation savings) and led the Boston School Committee to unanimously approve the first school start time reform in 30 years in December 2017.

However, the \gls{bps} experience also exposed the significant equity dimensions of school transportation: economically disadvantaged high school students were more likely to start before 7:30 AM. In fact, a 0.78 correlation was observed between median household income and elementary school start time. The attempted start-time reform, which would have reduced the proportion of high school students starting before 8:00 AM from 74\% to 6\%, was ultimately not implemented due to the magnitude of change required (85\% of schools would have had their times altered) and insufficient community engagement with those affected. Thus, this case powerfully illustrates both the potential and the institutional challenges of optimization-driven transit policy.

\subsection{Bounded Formulations for School Bus Scheduling}

\textcite{zeng2022Bounded} introduced a new time-indexed \gls{ilp} formulation for the school bus scheduling problem that jointly determines school start times and bus route operation times to minimize transportation cost. Their key contribution is a strengthened formulation obtained by characterizing the convex hulls of variable subspaces, exploiting the semi-decomposable block structure of the problem. Based on this strengthened \gls{lp} relaxation, they developed a dependent randomized rounding algorithm that yields near-optimal solutions for large-scale instances with a provable performance guarantee: the optimality gap is bounded by $O\left(\sqrt{R_{\text{max}}\log{T}\cdot OPT}\right)$, where $R_{\text{max}}$ is the maximum number of routes per school and $T$ is the number of time periods.

The work also introduced a robust version of the \gls{sbrp} that accounts for uncertainty in route sets across multiple scenarios which is critical, since school start times cannot be redesigned frequently but bus routes may vary from year to year. The generalization maintains similar theoretical guarantees while providing stakeholders with scheduling plans that remain consistent over time. This approach demonstrates how rigorous mathematical formulations can support strategic decision-making in school districts, providing tools with provable bounds rather than purely heuristic solutions.

\section{Equity in Transportation Planning}
\subsection{Conceptual Foundations of Transport Justice}

The question of how to evaluate and promote equity in transportation has evolved from a peripheral concern to a central challenge. \textcite{martensTransportJustice2016} developed a comprehensive philosophical framework for transport justice, arguing that governments have a fundamental duty to provide virtually every person with adequate transportation. Drawing on the work of Rawls and Dworkin, \textcite{martensTransportJustice2016} proposed that the central concern should be ``accessibility,'' the ease with which people can reach out-of-home activities, and developed the concept of a ``sufficiency threshold,'' below which accessibility levels constitute an injustice.

\textcite{martensTransportJustice2016} also introduced the \gls{afi}, adapted from the Foster-Greer-Thorbecke income poverty measure, to quantify the severity of accessibility deficiency in a region. The \gls{afi} captures both the prevalence and intensity of accessibility shortfalls across population groups, which allows planners to identify and prioritize groups most affected by inadequate transportation access. Crucially, \textcite{martensTransportJustice2016} argued that the measurement of accessibility poverty is synonymous with the measurement of the fairness of a transportation system. This is a stronger claim than in the income domain, justified by the inevitability of spatial disparities in accessibility and the bounded nature of time as a resource.

\subsection{Measuring Transportation Equity: Methods and Limitations}

Extending this work, \textcite{vanweeEvaluatingTransportEquity2021} provided a systematic review of how equity is used for evaluating transport policy impacts, covering metrics including the Gini index, Suits index, Palma index, Theil index, and Atkinson index. Their analysis confirmed that accessibility is the most commonly evaluated topic and the Gini index is by far the most frequently used metric, likely because of its ease of interpretability and communication. The review also identified significant research gaps: studies on distributions of safety and environmental effects using quantitative metrics are scarce, and the literature on citizens' preferences regarding equitable distributions of transport impacts is ``in its infancy.''

Here, \textcite{vanweeEvaluatingTransportEquity2021} distinguished several equity types relevant to transportation, such as horizontal, vertical, territorial, and spatial equity, and linked these to the domains of accessibility, safety, and the environment. They also noted the importance of the World Health Organization's definition of equity as ``the absence of avoidable or remediable differences among groups of people,'' emphasizing that the qualifier ``avoidable or remediable'' is crucial for transport-related distributive justice, since some accessibility differences are unavoidable consequences of different spatial arrangements.

\textcite{karnerAdvancesPitfallsMeasuring2025} advanced this perspective by identifying fundamental limitations of the two most widely used approaches for evaluating transport justice: Gini coefficients/Lorenz curves for measuring accessibility inequality and needs-gap/transit desert approaches for measuring accessibility poverty. They identified five major limitations of Gini/Lorenz methods: emphasis on equality rather than equity, failure to capture vertical equity, unordered comparisons that ignore socioeconomic characteristics, ambiguous direction of inequality, and reporting only within-group rather than between-group inequalities.

As alternatives, \textcite{karnerAdvancesPitfallsMeasuring2025} proposed the concentration index, which orders populations by socioeconomic status and can indicate both the direction and strength of the relationship between accessibility and social position. They also advocated for threshold-based approaches aligned with sufficientarian concerns, where injustice is defined not by inequality but by the failure to provide adequate levels of accessibility. Their analysis of COVID-19-related transit service cuts in Washington, D.C. demonstrated that standard Gini/Lorenz methods could ``easily mislead analysts, decision makers, and the public'' about equity impacts.

\subsection{Equity in Bus Network Redesigns}

\textcite{situTitleVIEnough2025} conducted a comprehensive review of equity analyses associated with bus network redesign projects across the United States. Their study found that most agencies rely on systemwide service equity analyses using U.S. Census data, which can obscure localized disparities, while the few agencies conducting route-by-route analyses were more effective at identifying adverse impacts on Title VI and \gls{ej} populations. The study proposed an enhanced stepwise framework for service equity analysis that extends the \gls{fta}'s standard four-level framework with five critical additions: use of rider-based data, context-sensitive measures, route-level analysis, recognition of non-Title VI and non-\gls{ej} populations (including seniors and people with disabilities), and integration of lived experience.

A recurring finding was the contrast between patronage-oriented service (high-frequency routes in high-demand corridors) and coverage-oriented service (broader geographic coverage to serve transit-dependent populations), with agencies typically prioritizing patronage at the risk of exacerbating existing inequalities. The study also revealed that local news reports consistently surfaced equity concerns that agency analyses had failed to detect: limited service for transit-dependent riders, longer walking distances, and challenges for older adults.

\subsection{Computational Approaches to Equitable Transit Design}

\textcite{paviaDesigningEquitableTransit2023} developed a mathematical formulation for transit network design that explicitly incorporates different notions of equity, welfare, and priority through a \gls{milp} based on a piecewise linear utility function. Their framework operationalizes three social welfare functions (utilitarian, Rawlsian max-min, and leximax) and demonstrates how explicitly accounting for diverse priority scores among origin-destination pairs affects network design outcomes.

Their experimental results on Chattanooga, Tennessee revealed several key insights: at moderate budgets, explicitly modeling priorities can significantly benefit groups with the highest transit need; the Rawlsian formulation ensures all regions are served but may reduce average utility compared to the utilitarian approach; and the leximax formulation can iteratively improve worst-case utility at the cost of average utility. These findings suggest that cities operating within moderate budget ranges can meaningfully improve equity outcomes through priority-aware optimization, though the choice between how to conceptualize equity remains a normative decision for policymakers.

\section{Applied Category Theory and the Co-Design Framework}
\subsection{Monotone Co-Design Theory}

The mathematical foundation for the approach this thesis will be taking draws on the monotone theory of co-design developed by \textcite{censi2016Mathematical}, which provides a formal methodology for designing complex systems composed of interconnected subsystems/components in a modular and compositional fashion. The theory's fundamental building block is the \gls{mdpi}, defined as a tuple consisting of a set of implementations, a functionality space, and a resource space, connected by monotone provision and requirement maps.

Individual \glspl{mdpi} can be connected through series, parallel, and loop composition to form \glspl{cdpi}, which allow for decomposing large problems into manageable subproblems. The key theoretical result is that for any given \gls{mcdp} where functionality and resources are complete partial orders and the feasibility relation is both monotone and Scott continuous, the minimal resources required to provide a given functionality can be characterized as a least fixed point. This approach handles multi-objective, nonconvex, nondifferentiable, and noncontinuous optimization problems defined over potentially non-continuous spaces, which are all properties that are characteristic of real-world transit design problems.

The compositionality property is especially powerful: because the composition of monotone \glspl{mdpi} is itself a monotone \gls{mdpi}, hierarchical problems can be decomposed without a loss of formal guarantees. Rather than requiring standard assumptions of linearity, continuity, or convexity, this method asks only for general monotonicity assumptions (for example, that the cost of serving a bus route does not decrease as the route becomes longer).

\subsection{Application to Mobility Systems}

\textcite{zardiniCoDesignEnableUserFriendly2023} demonstrated the application of monotone co-design theory to future mobility system design, developing a framework for co-designing systems involving \glspl{av}, micromobility solutions, and public transit. This framework models transportation systems as an edge-labeled directed graph composed of road, public transportation, and walking layers, with travel demands represented as multi-commodity flows.

Individual components (autonomous vehicle design, micromobility vehicle design, public transit scheduling) are formalized as separate \glspl{mdpi} and composed to capture systemic interdependencies. For example, the \gls{av} \gls{mdpi} selects the labeled graph on which the \gls{av} provider wants to operate based on the achievable speed, which in turn determines both the serviced network and operational costs. The full co-design problem computes the Pareto front of solutions, trading off travel time, transportation costs, and externalities. This then provides actionable information to stakeholders in the mobility ecosystem.

This method is applied to a case study of Washington, D.C. using \gls{osm} road network data, \gls{gtfs} public transit data, and travel demand datasets. Here, the co-design framework showcased how different investment scenarios produce different rational solutions along the Pareto front. The modularity of this approach also allows stakeholders to focus on their individual \glspl{mdpi}, while the compositional framework ensures overall system-level coherence.

\section{Data Infrastructure for Transit Planning}
\subsection{\glsentryfull{gtfs}}

\Gls{gtfs} \cite{Overview} has become the standard open format for public transportation schedule and geographic information. A \gls{gtfs} feed consists of a collection of text files, including \verb!agency!, \verb!routes!, \verb!trips!, \verb!stops!, \verb!stop_times!, and \verb!calendar! data, that together describe a transit system's services, schedules, and spatial configuration. The specification supports both static schedule data (\gls{gtfs} Schedule) and real-time updates (\gls{gtfs} Realtime) covering trip delays, service alerts, and vehicle positions.

\Gls{gtfs} data has found wide application beyond its original use in consumer-facing trip planners, extending to multi-modal trip planning, accessibility analysis, timetable creation, planning and analysis tools, and public information displays. The \gls{fta}'s STOPS tool uses \gls{gtfs} data to represent transit services for project evaluation, enabling representation of both existing services and proposed modifications. The standardized nature of \gls{gtfs} feeds enables systematic comparison across transit agencies and provides a ready-made data source for optimization frameworks. \cite{nchrp08-119General, 2019STOPS}

\subsection{OpenStreetMap and Road Network Data}

\Gls{osm} provides detailed, openly available street network data that includes information for routing by car, foot, bicycle, and other modes. Data available through \gls{osm} has been increasingly used in transportation planning applications, from pedestrian route planning to transit accessibility analysis. As previously mentioned, \textcite{zardiniCoDesignEnableUserFriendly2023} leveraged \gls{osm} data for their mobility co-design framework, deriving road arc lengths, speed limits, and capacity constraints from the open-source network. The crowdsourced nature of \gls{osm} data means coverage and quality varies by location, but for many urban areas, the data is sufficiently detailed to support transit network optimization. \cite{openstreetmapwikicontributors2026Routing, 2024Complexity}

\subsection{Census and American Community Survey Data}

U.S. Census and \gls{acs} data provide the demographic and socioeconomic information essential for equity-aware transit planning. \textcite{bertsimasOptimizingSchoolsStart2019} used \gls{acs} data on median household income to reveal the correlation between neighborhood affluence and favorable school start times in Boston, while \textcite{karnerAdvancesPitfallsMeasuring2025} used demographic data to construct the socioeconomic ordering necessary for concentration index calculations. \textcite{paviaDesigningEquitableTransit2023} used Census data on income, commuting characteristics, and origin-destination employment statistics to construct priority scores for transit network design. The integration of Census/\gls{acs} data with transit network data enables the kind of distributional analysis that equity-aware planning demands.

\section{Synthesis and Research Gaps}

This literature review reveals several critical gaps that I hope to explore. First, while substantial progress has been made in both bus network optimization and transportation equity measurement, these domains remain largely disconnected in practice. Optimization models for transit network design typically treat equity as an afterthought (either as an after-the-fact evaluation criterion or as a simplified constraint) rather than as a native component of the design process.

Second, existing optimization frameworks for transit design tend to be problem-specific and difficult to adapt to different planning contexts. The successful \gls{bps} implementation required substantial customization to handle Boston-specific constraints, and the authors noted that ``problem details often vary between districts.'' There is a need for general-purpose frameworks that can accommodate the varying requirements of different planners and contexts while maintaining formal guarantees on solution quality.

Third, the data integration challenge in transit planning remains formidable. Transit network design requires synthesizing road network data (from \gls{osm}), transit service data (from \gls{gtfs}), travel demand data (from \gls{od} surveys), and demographic data (from Census/\gls{acs}), each with different spatial and temporal resolutions, data models, and quality characteristics. While individual data sources are increasingly available, the tools for merging them into coherent planning inputs remain provisional.

Fourth, the monotone co-design framework from applied category theory offers a promising but underexplored avenue for addressing these integration challenges. \textcite{zardiniCoDesignEnableUserFriendly2023} demonstrated its applicability to mobility system design, but this application focused on autonomous vehicles and micromobility rather than bus network redesigns, and did not incorporate equity considerations. The categorical framework's emphasis on compositionality, monotonicity, and Pareto-optimal solutions aligns well with the multi-objective, multi-stakeholder nature of equitable transit design, but this connection has not been explored or formalized in the literature; this thesis takes a first step toward that formalization (see \autoref{ch:conclusions}).

This thesis addresses these gaps by developing a co-design framework that uses the mathematical infrastructure of applied category theory to create a compositional, data-integrated approach to bus network redesign. By formalizing datasets, stakeholder objectives, and optimization problems as interconnected \glspl{mdpi} within a co-design problem, this framework provides a method for incorporating equity alongside cost and service quality in the design process. The case study of school bus transit in Framingham, Massachusetts provides a concrete demonstration of the framework's applicability and efficacy in a real-world planning context, building on the rich literature on school bus routing while extending it with equity-aware and novel mathematical methods.